\documentclass[12pt,a4paper]{article}
\usepackage[utf8]{inputenc}
\usepackage[spanish,es-noquoting]{babel}
\usepackage[T1]{fontenc}
\usepackage{lmodern}
\usepackage{geometry}
\geometry{top=2cm, bottom=2cm, left=2.5cm, right=2.5cm}
\usepackage{xcolor}
\usepackage{listings}
\usepackage{enumitem}
\usepackage{fancyhdr}
\usepackage{hyperref}
\usepackage{tikz}
\usetikzlibrary{arrows.meta, positioning, calc, shapes}

\definecolor{codegreen}{rgb}{0,0.6,0}
\definecolor{backcolour}{rgb}{0.95,0.95,0.92}

\lstset{
    language=C,
    backgroundcolor=\color{backcolour},
    commentstyle=\color{codegreen},
    keywordstyle=\color{blue},
    basicstyle=\ttfamily\small,
    breaklines=true,
    frame=single,
    numbers=left,
    tabsize=4,
    extendedchars=true,
    showstringspaces=false,
    literate={á}{{\'a}}1 {é}{{\'e}}1 {í}{{\'i}}1 {ó}{{\'o}}1 {ú}{{\'u}}1 {ñ}{{\~n}}1
}

\pagestyle{fancy}
\fancyhf{}
\lhead{Monitorías Sistemas Operativos 2026-I}
\rhead{Pre-Taller: Tuberías}
\cfoot{\thepage}

\title{\textbf{Pre-Taller: El Anillo de Cifrado Distribuido}}
\author{Malak Sanchez \\ \small Monitorías de Sistemas Operativos}
\date{\today}

\begin{document}

\maketitle

\section*{Introducción}
Este ejercicio pone a prueba la capacidad de sincronización de flujos continuos de datos utilizando tuberías POSIX (\texttt{pipes}). Se busca implementar un sistema modular donde cada proceso cumpla una función de transformación específica.

\section*{Problema 1: Pipeline de Procesamiento}
Diseñe un programa que implemente un circuito de cifrado en etapas. El proceso padre enviará un mensaje de texto plano a través de una primera tubería hacia su primer hijo. Este primer hijo (Etapa 1) aplicará un algoritmo de \textbf{Cifrado César} desplazando 3 posiciones las letras y enviará el resultado temporal al segundo hijo (Etapa 2), el cual desplazará las letras 2 posiciones adicionales antes de devolver la cadena cifrada final al padre.

\begin{center}
\begin{tikzpicture}[node distance=1.5cm and 2.5cm, font=\small, >=Stealth]
    \node[draw, rectangle, fill=blue!10, minimum width=3cm] (padre) {PADRE};
    
    \node[draw, circle, fill=orange!20, right=of padre] (h1) {Etapa 1 (+3)};
    \node[draw, circle, fill=orange!20, below=of h1] (h2) {Etapa 2 (+2)};
    
    \draw[->, thick, color=blue!70!black] (padre) -- (h1) node[midway, above] {Pipe A};
    \draw[->, thick, color=red!70!black] (h1) -- (h2) node[midway, right] {Pipe B};
    \draw[->, thick, color=green!60!black] (h2) -| (padre) node[pos=0.25, below] {Pipe C};
\end{tikzpicture}
\end{center}

\textbf{Ejemplo de ejecución:} 
\begin{lstlisting}[numbers=none]
$ ./anillo_cifrado "HOLA"
[Padre] Mensaje original: HOLA
[Hijo 1] Cifrando +3... (KROD)
[Hijo 2] Cifrando +2... (MTPF)
[Padre] Mensaje cifrado final: MTPF
\end{lstlisting}

\section*{Requerimientos}
\begin{enumerate}
    \item Establecer una jerarquía de procesos y un flujo unidireccional estricto entre ellos.
    \item Cerrar de forma segura todos los descriptores de lectura/escritura sobrantes en cada proceso.
    \item Implementar correctamente la rotación de caracteres para el cifrado solicitado.
\end{enumerate}

\section*{Criterios de Evaluación}
\begin{itemize}
    \item \textbf{Uso de IPC:} Manejo experto de tuberías y descriptores de archivo.
    \item \textbf{Lógica:} Implementación correcta del algoritmo de sustitución.
    \item \textbf{Robustez:} Limpieza de recursos y finalización controlada de procesos.
\end{itemize}

\end{document}
